% Prajwal Khanal — Curriculum Vitae
%
% Plain LaTeX, no CV class. The layout is the classic academic one-pager:
% small-caps section labels in a left column, content in a right column,
% dates set right-aligned in italic.
%
% Build:  pdflatex cv.tex     (run it once; there is no bibliography)
%
% The face is Times, to match the website. Delete the mathptmx line below to
% get Computer Modern, the default LaTeX look.

\documentclass[11pt,letterpaper]{article}

\usepackage[utf8]{inputenc}
\usepackage[T1]{fontenc}
\usepackage{mathptmx}                 % Times. Remove for Computer Modern.
\usepackage[margin=0.7in]{geometry}
\usepackage{enumitem}
\usepackage[dvipsnames]{xcolor}
\usepackage[hidelinks]{hyperref}

% The same blue as the website and the existing CV.
\definecolor{cvblue}{RGB}{0,112,192}
\hypersetup{colorlinks=true, urlcolor=cvblue, linkcolor=cvblue}

\pagestyle{empty}
\setlength{\parindent}{0pt}
\setlength{\parskip}{0pt}

% ---------------------------------------------------------------------------
% Section: a small-caps label in the left margin, content to the right.
% ---------------------------------------------------------------------------
\newlength{\cvlabelwidth}
\setlength{\cvlabelwidth}{1.05in}

\newenvironment{cvsection}[1]{%
  \begin{list}{}{%
    \setlength{\leftmargin}{\cvlabelwidth}%
    \addtolength{\leftmargin}{0.15in}%
    \setlength{\labelwidth}{\cvlabelwidth}%
    \setlength{\labelsep}{0.15in}%
    \setlength{\itemindent}{0pt}%
    \setlength{\listparindent}{0pt}%
    \setlength{\topsep}{0.10in}%
    \setlength{\parsep}{0pt}%
    \setlength{\itemsep}{0pt}%
    % A parbox, so a label written with \\ breaks over two lines instead of
    % running together.
    % \smash zeroes the box's height and depth. Without it a two-line label
    % ("Contact \\ Information") makes the item's first line box as deep as
    % the label, which shows up as a gap between the first and second lines
    % of the content beside it.
    \renewcommand{\makelabel}[1]{%
      \smash{\parbox[t]{\cvlabelwidth}{\raggedright\textsc{##1}}}}%
  }%
  \item[#1]%
}{\end{list}}

% One entry: bold title and detail on the left, italic dates on the right.
% Facing minipages rather than \hfill, so a long title wraps within its own
% column instead of shoving the date onto the following line.
\newcommand{\entry}[3]{%
  \noindent
  \begin{minipage}[t]{0.70\linewidth}\raggedright\textbf{#1}#2\end{minipage}%
  \hfill
  \begin{minipage}[t]{0.28\linewidth}\raggedleft\textit{#3}\end{minipage}%
  \par
}

% Same, but with the title in roman: for awards and teaching, where there is
% no institution to set apart.
\newcommand{\plainentry}[2]{%
  \noindent
  \begin{minipage}[t]{0.70\linewidth}\raggedright #1\end{minipage}%
  \hfill
  \begin{minipage}[t]{0.28\linewidth}\raggedleft\textit{#2}\end{minipage}%
  \par
}

% A publication. The vertical space between them lives here, in one place.
\newcommand{\pub}[1]{\vspace{0.045in}#1\par}

% A heading inside a section, e.g. "Peer-reviewed".
\newcommand{\subhead}[1]{\vspace{0.07in}\textit{#1}\par\vspace{-0.02in}}

% A line of detail under an entry.
\newcommand{\detail}[1]{%
  \begin{list}{}{%
    \setlength{\leftmargin}{1.1em}\setlength{\topsep}{0pt}%
    \setlength{\parsep}{0pt}\setlength{\itemsep}{0pt}%
    \setlength{\partopsep}{0pt}}%
  \item #1
  \end{list}%
}

% Space between entries inside a section.
\newcommand{\gap}{\vspace{0.06in}}

% ---------------------------------------------------------------------------

\begin{document}

\begin{center}
  {\large\textbf{Prajwal Khanal \textemdash{} Curriculum Vitae}}
\end{center}
\vspace{-0.08in}
\rule{\textwidth}{0.8pt}

\begin{cvsection}{Contact\\Information}
  ITC Faculty, University of Twente \hfill \href{mailto:p.khanal@utwente.nl}{p.khanal@utwente.nl}\par
  Enschede, the Netherlands \hfill \href{mailto:ktm.prajwalkhanal@gmail.com}{ktm.prajwalkhanal@gmail.com}\par
  \href{https://prajwalkhanal.earth}{prajwalkhanal.earth} \hfill
  \href{https://github.com/prajzwal08}{github.com/prajzwal08} $\cdot$
  \href{https://scholar.google.com/citations?user=UtzNqTwAAAAJ}{Google Scholar}\par
\end{cvsection}

\begin{cvsection}{Research\\Interests}
  Soil--vegetation--atmosphere interaction; water and carbon fluxes; land surface
  modelling and deep learning, and where the two can inform each other.
\end{cvsection}

\begin{cvsection}{Education}
  \entry{PhD, Geo-Information Science and Earth Observation}{}{Feb 2023---\ldots}
  \detail{University of Twente, the Netherlands}
  \detail{Modelling soil--vegetation--atmosphere interaction, in both process-based
          land surface modelling and deep learning approaches to water and carbon fluxes.}
  \gap
  \entry{M.Sc., Environmental Engineering,}{ Technical University of Munich}{Nov 2020---Dec 2022}
  \detail{Graduated 1.4/5.0 on the German scale, where 1.0 is the highest mark.}
  \detail{Thesis on near-surface versus sub-surface soil moisture in vegetation
          functioning, published in \textit{Biogeosciences} (2024).}
  \gap
  \entry{B.E., Civil Engineering,}{ Institute of Engineering, Tribhuvan University, Nepal}{Nov 2014---Aug 2018}
  \detail{Ranked 20th of 192, 77.7\%. Thesis on flood inundation mapping of the Babai River.}
\end{cvsection}

\begin{cvsection}{Employment}
  \entry{PhD Researcher,}{ ITC, University of Twente}{Feb 2023---\ldots}
  \detail{WUNDER: ecohydrological responses of agricultural and nature ecosystems
          in the Netherlands to water use and drought.}
  \detail{Producing field-scale soil moisture and fluxes for agricultural drought monitoring.}
  \gap
  \entry{Research Intern,}{ Max Planck Institute for Biogeochemistry, Jena}{Nov 2020---Feb 2022}
  \detail{Quantified the relative control of near-surface soil moisture versus terrestrial
          water storage on global vegetation functioning, from satellite observation alone.}
  \gap
  \entry{Student Research Assistant,}{ Technical University of Munich}{Nov 2021---Mar 2022}
  \detail{Flood inundation maps for the Ouémé River, Benin (FURIFLOOD) and the Naryn
          River, Kyrgyzstan (ÖkoFlussPlan), using HEC-RAS 1D unsteady-flow models.}
  \gap
  \entry{Working Student,}{ Jena-Geos Ingenieurbüro GmbH, Munich}{May 2021---Dec 2021}
  \detail{Excel and VBA tool for greenhouse-gas accounting across stationary combustion,
          refrigeration and transport.}
  \gap
  \entry{Civil Engineer (freelance),}{ ECEP, Kathmandu, Nepal}{Dec 2018---Dec 2022}
  \detail{Business plans for community-based water supply projects, for the Department
          of Water Supply and Sanitation, Government of Nepal.}
  \detail{Smart water supply master plan for the development of Lumbini as a smart city,
          for the Department of Urban Development and Building Construction.}
\end{cvsection}

\begin{cvsection}{Publications}
  \subhead{Peer-reviewed}
  \pub{Moutahir, H., Sulzer, M., Kiese, R., \ldots{} \textbf{Khanal, P.}, et al. (2026).
  Matching scales of eddy covariance measurements and process-based modelling:
  assessing spatiotemporal dynamics of carbon and water fluxes in a mixed forest in
  Southern Germany. \textit{Biogeosciences}, 23, 1719--1738.
  \href{https://doi.org/10.5194/bg-23-1719-2026}{doi}}

  \pub{Zeng, Y., Alidoost, F., Schilperoort, B., \ldots{} \textbf{Khanal, P.}, et al. (2025).
  Towards an open soil-plant digital twin based on the STEMMUS-SCOPE model, following
  open science. \textit{Computers \& Geosciences}, 205, 106013.
  \href{https://doi.org/10.1016/j.cageo.2025.106013}{doi}}

  \pub{\textbf{Khanal, P.}, Hoek van Dijke, A. J., Schaffhauser, T., et al. (2024).
  Relevance of near-surface soil moisture vs.\ terrestrial water storage for global
  vegetation functioning. \textit{Biogeosciences}, 21(6), 1533--1547.
  \href{https://doi.org/10.5194/bg-21-1533-2024}{doi}}

  \pub{Bhattarai, P., \textbf{Khanal, P.}, Tiwari, P., et al. (2019).
  Flood inundation mapping of Babai basin using HEC-RAS and GIS.
  \textit{Journal of the Institute of Engineering}, 15(2), 32--44.
  \href{https://doi.org/10.3126/jie.v15i2.27639}{doi}}

  \subhead{Under review}
  \pub{\textbf{Khanal, P.}, Zeng, Y., Daoud, M. G., Beckers, J., Su, Z.
  Investigating the impact of multilayer resistance schemes on simulating land surface
  temperature and turbulent fluxes: evaluation within the STEMMUS-SCOPE model.
  \href{https://papers.ssrn.com/sol3/papers.cfm?abstract_id=6389498}{Preprint on SSRN}.}

  \pub{\textbf{Khanal, P.}, et al.
  Deep learning soil moisture and fluxes with multimodal foundation-model embeddings.}

  \subhead{Conference contributions}
  \pub{\textbf{Khanal, P.}, et al. (2026). Bridging the scale gap: daily 10-metre soil
  moisture, evapotranspiration and carbon fluxes via multimodal deep learning.
  \textit{Workshop on Machine Learning for Land and Hydrology}, ECMWF, Reading, UK.}

  \pub{\textbf{Khanal, P.}, Hoek van Dijke, A., Zeng, Y., Orth, R. (2023). Near-surface vs.\
  sub-surface soil moisture impacts on vegetation functioning.
  \textit{EGU General Assembly}, Vienna, EGU23-7670.
  \href{https://doi.org/10.5194/egusphere-egu23-7670}{doi}}

  \pub{\textbf{Khanal, P.} (2019). Flood inundation mapping: an inference on the outbreak of
  flood-borne diseases. \textit{International Conference on Health Engineering in Disaster}.}
\end{cvsection}

\begin{cvsection}{Awards}
  \plainentry{Second place, Hackathon on Machine Learning for Earth System Modelling}{Aug 2026}
  \detail{Quantization in foundation models. Supported by CESOC, ECMWF,
          TRA Modelling (Bonn), FZ Jülich and the University of Cologne.}
  \gap
  \plainentry{Deutschlandstipendium, national merit scholarship, TU Munich}{2021/22}
  \gap
  \plainentry{Institute of Engineering merit-based scholarship, Tribhuvan University}{2014---2018}
\end{cvsection}

\begin{cvsection}{Teaching}
  \plainentry{Teaching Assistant, Technical University of Munich}{2021---2022}
  \detail{Hydrological and Environmental River Basin Modelling; Integrated Water
          Resources Management; Rapidly Varying Flow.}
  \gap
  \plainentry{Assistant Lecturer, Cosmos College of Management and Technology, Nepal}{Nov 2018---Jan 2021}
  \detail{Designed and taught an elective on Python and GIS for water resources;
          tutorials in Fluid Mechanics and Hydraulics.}
\end{cvsection}

\begin{cvsection}{Software}
  \entry{WUNDER data viewer}{, ITC, University of Twente}{2026}
  \detail{Interactive app for the WUNDER weather and soil moisture loggers, so growers
          at the food forests can follow drought as it develops.
          \href{https://wunder-project.streamlit.app}{App} $\cdot$
          \href{https://github.com/prajzwal08/wunder-project}{Source}}
\end{cvsection}

% CHECK THIS SECTION. It is assembled from the work described above rather than
% from a list you wrote, so the groupings are a guess. Correct it before sending.
\begin{cvsection}{Skills}
  \textit{Modelling:} STEMMUS-SCOPE, land surface modelling, HEC-RAS 1D unsteady flow\par
  \textit{Programming:} Python, deep learning frameworks, GIS, Excel/VBA\par
  \textit{Data:} eddy covariance and flux-tower data, satellite soil moisture, ICOS\par
  \textit{Languages:} English, Nepali\par
\end{cvsection}

\end{document}
